\clearpage
\begin{explanationbox}
\textbf{第一眼：先看最简单的三个点。}

取 $C:y^2=x$、$M(2,0)$、$r=1$，从抛物线顶点 $A_1=O(0,0)$ 作两条圆的切线。它们与抛物线再次相交于
\[
A_2=(3,\sqrt3),\qquad A_3=(3,-\sqrt3).
\]
$A_1A_2:\ x-\sqrt3y=0$ 到 $M$ 的距离为 $2/\sqrt{1+3}=1$；$A_1A_3$ 同理。第三条弦\textbf{恰好竖直}：
\[
A_2A_3:\ x=3,\qquad d(M,A_2A_3)=|3-2|=1.
\]
\textbf{这就是三条切线闭合的第一幅画面。}

\medskip
\textbf{第二眼：把公共顶点移走，第三条仍然碰圆一下。}
\begin{center}
% Equal scale across both panels; both panels use exactly the same circle.
% second A=(1/4,1/2), B,C have t=(2+-sqrt37)/3.
\begin{tikzpicture}[x=0.62cm,y=0.62cm,>=Stealth,font=\scriptsize]
\begin{scope}
\draw[gray!55,->] (-0.45,0)--(8.25,0) node[right] {$x$};
\draw[gray!55,->] (0,-3.10)--(0,3.15) node[above] {$y$};
\draw[blue!70!black,very thick,domain=-2.75:2.75,samples=130,variable=\t] plot ({\t*\t},{\t});
\draw[teal!75!black,thick,fill=teal!9] (2,0) circle (1);
\coordinate (A) at (0,0);\coordinate (B) at (3,{sqrt(3)});\coordinate (C) at (3,{-sqrt(3)});
\draw[orange!80!black,very thick] (A)--(B)--(C)--cycle;
\foreach \p in {A,B,C} \fill[orange!80!black] (\p) circle (2.1pt);
\node[below left] at (A) {$A$};\node[above left] at (B) {$B$};\node[below left] at (C) {$C$};
\fill[teal!80!black] (2,0) circle(1.7pt);\node[below] at (2,0) {$M$};
\node[font=\normalsize] at (4,3.55) {顶点 $A=(0,0)$};
\end{scope}
\begin{scope}[xshift=6.5cm]
\draw[gray!55,->] (-0.45,0)--(8.25,0) node[right] {$x$};
\draw[gray!55,->] (0,-3.10)--(0,3.15) node[above] {$y$};
\draw[blue!70!black,very thick,domain=-2.75:2.75,samples=130,variable=\t] plot ({\t*\t},{\t});
\draw[teal!75!black,thick,fill=teal!9] (2,0) circle (1);
\pgfmathsetmacro{\bb}{(2+sqrt(37))/3}
\pgfmathsetmacro{\cc}{(2-sqrt(37))/3}
\coordinate (A) at (0.25,0.5);\coordinate (B) at ({\bb*\bb},{\bb});\coordinate (C) at ({\cc*\cc},{\cc});
\draw[orange!80!black,very thick] (A)--(B)--(C)--cycle;
\foreach \p in {A,B,C} \fill[orange!80!black] (\p) circle (2.1pt);
\node[below left] at (A) {$A$};\node[above right] at (B) {$B$};\node[below right] at (C) {$C$};
\fill[teal!80!black] (2,0) circle(1.7pt);\node[below] at (2,0) {$M$};
\node[font=\normalsize] at (4,3.55) {顶点 $A=(\frac14,\frac12)$};
\end{scope}
\end{tikzpicture}

\end{center}
两图中的抛物线、圆和坐标尺度都没有变。左图从原点出发，右图把第一个顶点改为 $A_1=(1/4,1/2)$；另外两点位置也变了，但第三条弦仍与同一个圆相切。\textbf{变的是三个点；不变的是第三边相切。}

\medskip
\textbf{第三眼：为什么原题要找“和与积”？}

把 $A_i$ 写为 $P(y_i)=(y_i^2/k,y_i)$。第三条弦是
\[
\boxed{kx-(y_2+y_3)y+y_2y_3=0}.
\]
所以只需求\textbf{$y_2+y_3$ 与 $y_2y_3$}；第 03 节的和积图直接标出这两个量怎样确定第三条弦。

\medskip
\textbf{它替你省掉什么？} 遇到“同一个点连出两条切圆的弦，要求第三条弦”时，不必三次联立求交点。先看圆心是不是特别位置；是则直接判断相切，若要完整证明，再做和、积的硬算。
\end{explanationbox}
